% Conference Paper -- a compact two-column paper on US Letter.
% The page design (grid, XCharter and Lato, headings, title block, accent
% colour) lives in conference.sty; this file holds the content.
% Edit the title, author grid, abstract and keywords in the metadata block.
% Build with: latexmk -pdf main.tex
\documentclass[10pt, twocolumn]{article}
\usepackage{conference}

\usepackage{tikz}
\usetikzlibrary{arrows.meta, positioning, fit, backgrounds}
\usepackage{pgfplots}
\pgfplotsset{compat=1.18}
\usepgfplotslibrary{groupplots}
\usepackage{siunitx}
\sisetup{detect-weight, mode=text, group-minimum-digits=4}
\usepackage[ruled, vlined, linesnumbered, noend]{algorithm2e}
\SetAlFnt{\small}
\SetAlCapFnt{\footnotesize\sffamily}
\SetAlCapNameFnt{\footnotesize\sffamily}
\SetKwComment{tcp}{\textcolor{black!50}{// }}{}
\renewcommand{\CommentSty}[1]{\textcolor{black!55}{\itshape #1}}

% Numeric references in a small size, compressed into ranges.
\usepackage[backend=biber, style=numeric-comp, sorting=nyt, maxbibnames=3,
  giveninits=true, doi=false, url=false, isbn=false, eprint=false]{biblatex}
\addbibresource{refs.bib}
\renewcommand*{\bibfont}{\footnotesize}
\setlength{\bibitemsep}{1pt}

\usepackage[colorlinks, allcolors=accent]{hyperref}
\usepackage[capitalise, noabbrev]{cleveref}

% One look for every plot: grey baselines, the accent for our system.
\pgfplotsset{every axis/.append style={
  font=\sffamily\scriptsize, axis line style={black!45}, tick style={black!45},
  grid=major, grid style={black!8}, every axis plot/.append style={line width=.9pt},
  legend style={draw=none, fill=none, font=\sffamily\scriptsize}}}

% ---- Metadata --------------------------------------------------------------
\title{Lantern: Learned Prefetching for Tiered Key-Value Stores}
\venue{Proceedings of the Workshop on Data Systems, 2026}
\authorgrid{%
  \authorcell{Noor Haddad}{University of Toronto}{noor@cs.example.edu}\hfill
  \authorcell{Mateus Lacerda}{University of Toronto}{mlacerda@cs.example.edu}\hfill
  \authorcell{Yuki Tanabe}{Kioxia Research}{y.tanabe@example.com}\hfill
  \authorcell{Rachel Abernethy}{University of Toronto}{rabernethy@cs.example.edu}}
\ccsconcepts{Information systems $\rightarrow$ Key-value stores;
  Hierarchical storage management; Computing methodologies $\rightarrow$
  Supervised learning.}
\keywords{prefetching, tiered storage, key-value stores, learned systems}

\begin{abstract}
Key-value stores increasingly place hot data in DRAM and the long tail on
NVMe flash. Replacement policies decide what to keep in the fast tier, but
they act only after a miss, and the miss is what costs a flash read. We
present \sys{}, a prefetcher that predicts, for each key that is
accessed, the small set of keys likely to follow it within the next
millisecond and promotes them ahead of demand. A gradient-boosted model
over co-access features fits in 6\,MB per shard, is retrained online every
few minutes, and gates every prefetch by a cost-aware threshold so that
wrong guesses never evict useful data. On four production-derived traces
\sys{} lowers P99 read latency by 31--47\% over ARC and by 18\%
over the best sequence-model prefetcher, at under 3\% CPU overhead.
\end{abstract}

% The teaser spans both columns under the abstract on page one.
\teaser{%
  \begin{tikzpicture}[
      font=\sffamily\footnotesize, >={Stealth[length=4pt]},
      box/.style={draw=black!45, rounded corners=2pt, fill=white,
        minimum height=8.5mm, inner xsep=6pt, align=center},
      hot/.style={box, draw=accent, fill=accent!7, text=accent!80!black},
      tier/.style={draw=black!25, rounded corners=3pt, fill=black!3, inner sep=7pt}]
    \node[box] (client) {Client\\[-1pt]{\scriptsize get / put}};
    \node[box, right=9mm of client] (index) {Shard index\\[-1pt]{\scriptsize hash $\to$ tier}};
    \node[hot, right=9mm of index] (feat) {Co-access features\\[-1pt]{\scriptsize last $k$ keys, gaps}};
    \node[hot, right=9mm of feat] (model) {Gradient-boosted\\[-1pt]{\scriptsize successor model}};
    \node[hot, right=9mm of model] (gate) {Cost gate\\[-1pt]{\scriptsize $p\cdot c_{\mathrm{miss}} > c_{\mathrm{evict}}$}};
    \node[box, below=10mm of index, minimum width=34mm] (dram) {DRAM tier\\[-1pt]{\scriptsize 64 GB, LRU-K}};
    \node[box, right=14mm of dram, minimum width=34mm] (nvme) {NVMe tier\\[-1pt]{\scriptsize 2 TB, log-structured}};
    \begin{scope}[on background layer]
      \node[tier, fit=(dram)(nvme), label={[font=\sffamily\scriptsize, black!55]left:storage}] {};
      \node[tier, fill=accent!3, draw=accent!30, fit=(feat)(model)(gate),
        label={[font=\sffamily\scriptsize\bfseries, accent]above:\sys{}}] {};
    \end{scope}
    \draw[->, black!60] (client) -- (index);
    \draw[->, black!60] (index) -- (feat);
    \draw[->, accent] (feat) -- (model);
    \draw[->, accent] (model) -- (gate);
    \draw[->, black!60] (index) -- (dram);
    \draw[->, accent, thick] (gate.south) |- node[pos=.75, above, accent] {promote ahead of demand} (nvme.east);
    \draw[->, black!60] (nvme.west) -- node[above] {miss} (dram.east);
  \end{tikzpicture}\par\vspace{4pt}
  \captionof{figure}{\sys{} sits beside the shard index. Every access
  updates co-access features; the model proposes likely successors and a cost
  gate promotes only those whose expected saving exceeds the cost of the
  eviction they cause.}
  \label{fig:teaser}}

\begin{document}
\maketitle


\section{Introduction}

Large key-value deployments no longer fit in memory. Operators place a few
percent of the data in DRAM and the rest on NVMe flash, and rely on a
replacement policy to keep the hot set in the fast tier~\cite{varga2022tiered,
yang2020twitter}. Replacement has been studied for decades, from LRU-K to
ARC to learned approximations of Belady's
algorithm~\cite{oneil1993lruk,megiddo2003arc,song2020lrb}, and modern
policies capture most of what recency and frequency can explain.

What no replacement policy can do is avoid the first miss. A key that has
not been accessed recently will be read from flash, at roughly a hundred
times the latency of DRAM, and on production traces these cold reads
dominate the tail~\cite{yang2020twitter}. Yet many of them are predictable:
keys are accessed in groups that follow application logic, such as a user
record followed by its settings and its recent sessions. A prefetcher that
learned these groups could promote the followers while the leader is being
served.

Prefetching for storage is not new, and sequence models trained on access
logs predict the next key well in the
laboratory~\cite{hashemi2018,osei2023prefetch}. They are hard to deploy for
three reasons. They are large and slow to run per request; they drift as
applications change; and a wrong prefetch is not free, because it evicts a
key that might have been used. We address each in turn:

\begin{itemize}
  \item A compact \emph{successor model} built on co-access features and
    gradient-boosted trees~\cite{ke2019lightgbm}, small enough to run inline
    (\cref{sec:model}).
  \item \emph{Online retraining} on a sliding window, which tracks workload
    drift without an offline pipeline.
  \item A \emph{cost gate} that admits a prefetch only when its expected
    saving exceeds the expected cost of the eviction it causes
    (\cref{sec:gate}).
\end{itemize}

\noindent \Cref{sec:eval} evaluates \sys{} on four traces and
shows where it does not help.

\section{Background and Motivation}
\label{sec:background}

\subsection{Tiered key-value stores}
A tiered store keeps an index from keys to locations and serves reads from
DRAM when it can. On a miss it reads the value from flash, returns it, and
admits it to DRAM, evicting another key. Admission and eviction are the
policy's two levers; recent work tunes both, for example by filtering
one-hit wonders with ghost lists~\cite{brandao2024ghost}.

\subsection{Where the misses come from}
We analysed a week of requests to 32 shards of a social-graph store. After
ARC has warmed up, 71\% of the remaining DRAM misses are \emph{compulsory
within the window}: the key was not accessed in the preceding hour. Of
those, 58\% are accessed within one millisecond of another key with which
they were co-accessed at least five times in the previous day. This is the
opportunity \sys{} targets: misses that recency cannot predict but
co-access history can.

\paragraph{Why not a sequence model?}
A two-layer LSTM over the last 64 keys predicts the next key with 44\%
top-10 accuracy on our traces, but costs \SI{38}{\micro\second} per
inference on a server core, more than the DRAM read it is meant to protect.
Batched inference amortises the cost but delays the prediction past the
moment it is useful.

\subsection{Workload characterisation}
To understand which applications produce learnable successors, we grouped
the misses of the \emph{Graph} trace by the key prefix that identifies the
calling service. Three services account for 80\% of the predictable misses:
profile rendering, which reads a user record and then its settings and
avatar metadata; feed assembly, which reads a list of post identifiers and
then the posts; and notification delivery, which reads a device token after
each user record. In all three the successors are fixed by code rather
than by user behaviour, which is why they recur so reliably.

The remaining services issue reads whose successors depend on the user, for
example the next item in a search result. Their co-access counts are low and
their gaps long, so the model assigns them low probabilities and the gate
rarely admits them. This division explains a pattern visible throughout our
results: \sys{} is most effective on traces dominated by a few services
with fixed read sequences, and least effective on traces whose accesses are
driven by independent user choices.

Two properties of the traces matter for the design. First, the gap between
a leader and its successors is short: the median is \SI{140}{\micro\second}
and the 90th percentile is \SI{610}{\micro\second}, which is why the window
$\Delta$ is one millisecond and why a prediction is useless if it arrives
after a batching delay. Second, successor sets are small: 90\% of leaders
have at most twelve successors with a co-access count above five, so a
sketch of 32 entries loses little.

\subsection{Design goals}
These observations set four goals. First, a prediction must be ready before
the leader's read completes, which on DRAM is a few microseconds; anything
slower is a prediction about the past. Second, the model must adapt to
changes in application logic without an offline pipeline, because the
operators of a storage tier rarely control the applications above it.
Third, a wrong prediction must cost less than the benefit of a right one,
on every workload, including those with no structure at all. Fourth, the
memory spent on metadata must be small compared with the DRAM it protects,
or the same bytes would be better spent on a larger cache.

The first goal rules out sequence models on the request path, and the
second rules out models trained once and shipped. The third is the one most
prior prefetchers do not address: they are evaluated by accuracy, while the
quantity that matters is the change in tail latency once mistakes are paid
for. \sys{} is designed backwards from that quantity.

\section{Design}
\label{sec:model}

\sys{} runs inside each shard and observes every access. It has
three parts: a feature extractor, a successor model and a cost gate
(\cref{fig:teaser}).

\subsection{Co-access features}
For each key $k$ the extractor keeps a bounded \emph{successor sketch}
$S(k)$: the up to 32 keys that most often followed $k$ within a window
$\Delta$, with counts and the median gap between the two accesses. Sketches
are updated in constant time with a space-saving counter and are stored
next to the index entry, so a lookup costs no extra random access. For a
candidate pair $(k, j)$ with $j \in S(k)$ the features are the co-access
count and its decay-weighted variant, the median gap, the recency and
frequency of $j$, whether $j$ is already in DRAM, and the size of its value.

\subsection{Successor model}
A single gradient-boosted model per shard scores each candidate with the
probability $p(j \mid k)$ that $j$ is read within $\Delta$. We use 64 trees
of depth six, which take \SI{1.9}{\micro\second} to score 32 candidates and
occupy 6\,MB. The model is retrained every five minutes on labels collected
from the previous window, which are free: we observe whether each candidate
was in fact read. Retraining runs on a background core and swaps the model
atomically.

\subsection{Cost gate}
\label{sec:gate}

A prefetch is useful only if its expected benefit outweighs the harm it
does. Promoting $j$ saves a flash read with probability $p(j \mid k)$, but
it evicts the DRAM victim $v$, which may itself be read soon. Let
$c_{\mathrm{miss}}$ be the cost of a flash read and $q(v)$ the probability
that $v$ is read before it would have been evicted anyway, estimated from
its ARC list position. The gate admits $j$ when
\begin{equation}
\label{eq:gate}
p(j \mid k)\,c_{\mathrm{miss}} \;>\; q(v)\,c_{\mathrm{miss}} + c_{\mathrm{io}},
\end{equation}
where $c_{\mathrm{io}}$ charges the flash bandwidth the prefetch consumes.
Because both sides scale with $c_{\mathrm{miss}}$, the rule reduces to
$p(j \mid k) - q(v) > c_{\mathrm{io}} / c_{\mathrm{miss}}$, a threshold that
we set once per device. \Cref{alg:lantern} gives the complete access path.

\begin{algorithm}[t]
\caption{Access path with \sys{} prefetching.}
\label{alg:lantern}
\KwIn{accessed key $k$, shard index $I$, model $M$, threshold $\tau$}
serve $k$ from DRAM or flash and update $I$\;
update successor sketches of the last keys with $k$\;
$C \leftarrow S(k) \setminus \mathrm{DRAM}$ \tcp*{candidates not yet hot}
\ForEach{$j \in C$ in batches of 32}{
  $p_j \leftarrow M(\phi(k, j))$\;
}
\ForEach{$j \in C$ by decreasing $p_j$}{
  $v \leftarrow$ next eviction victim in DRAM\;
  \If{$p_j - q(v) > \tau$}{
    issue asynchronous read of $j$; promote on completion\;
  }
  \lElse{\textbf{break} \tcp*[f]{later candidates score lower}}
}
record $(k, C)$ as pending labels for retraining\;
\end{algorithm}

% Full-width floats appear one page after their source position.
\begin{figure*}[t]
  \centering
  \begin{tikzpicture}
    \begin{groupplot}[group style={group size=3 by 1, horizontal sep=14mm},
        width=0.34\textwidth, height=44mm]
      \nextgroupplot[title={Graph: read latency CDF}, xmode=log,
          xlabel={latency (\si{\micro\second})}, ylabel={fraction of reads},
          xmin=5, xmax=1000, ymin=0, ymax=1]
        \addplot[black!35, dashed] coordinates {(5,0) (8,.55) (12,.74) (40,.80) (100,.84) (200,.90) (377,.99) (800,1)};
        \addplot[black!60] coordinates {(5,0) (8,.58) (12,.80) (40,.85) (100,.89) (200,.95) (262,.99) (700,1)};
        \addplot[accent] coordinates {(5,0) (8,.61) (12,.86) (40,.90) (100,.94) (160,.98) (201,.99) (600,1)};
      \nextgroupplot[title={Precision and coverage vs.\ $\tau$},
          xlabel={gate threshold $\tau$}, ylabel={percent},
          xmin=0, xmax=0.5, ymin=0, ymax=100,
          legend style={at={(0.97,0.5)}, anchor=east}]
        \addplot[accent] coordinates {(0,38) (.05,52) (.1,63) (.15,73) (.2,79) (.3,86) (.4,90) (.5,92)};
        \addplot[black!60] coordinates {(0,81) (.05,74) (.1,66) (.15,58) (.2,49) (.3,34) (.4,22) (.5,14)};
        \legend{precision, coverage}
      \nextgroupplot[title={Recovery after a layout change},
          xlabel={minutes after change}, ylabel={precision (\%)},
          xmin=0, xmax=60, ymin=0, ymax=100,
          legend style={at={(0.97,0.03)}, anchor=south east}]
        \addplot[black!45, dashed] coordinates {(0,72) (5,40) (10,31) (20,30) (40,29) (60,29)};
        \addplot[accent] coordinates {(0,72) (5,41) (10,58) (15,69) (20,71) (40,72) (60,73)};
        \legend{frozen model, online}
    \end{groupplot}
  \end{tikzpicture}
  \caption{Left: CDF of read latency on \emph{Graph} for ARC (dashed), ARC
  with an LSTM prefetcher (grey) and Lantern. Centre: the gate threshold
  trades coverage of misses for precision of prefetches; the device-derived
  value is $\tau = 0.15$. Right: precision after the \emph{Shop} catalogue
  changes layout at minute zero, with a frozen model and with five-minute
  online retraining.}
  \label{fig:behaviour}
\end{figure*}

\section{Implementation}
\label{sec:impl}

We built \sys{} as a library of about 3{,}400 lines of C++ linked into the
storage server. Sketches live in a slab next to each hash-index bucket, so
updating them touches a cache line that the lookup has already loaded. The
model is compiled to a flat array of split thresholds and evaluated with
branch-free code, which is what keeps scoring under two microseconds; a
generic tree library was six times slower on the same model.

Prefetches are issued through the same asynchronous I/O queue as demand
reads but at a lower priority, so a burst of prefetches can never delay a
client read. A promoted key enters DRAM at the middle of the ARC recency
list rather than at its head: if the prediction was wrong the key leaves
quickly, and if it was right the first real access moves it to the head as
usual. Labels for retraining are written to a ring buffer of 4 million
entries per shard, and a background thread fits the next model with the
same feature code that serves requests, which removes a whole class of
training--serving skew.

\paragraph{Failure handling}
If retraining fails or produces a model whose validation precision falls
below the current one, the shard keeps the current model and logs the
event. If the I/O queue exceeds a depth limit, the gate raises its
threshold until the queue drains, so prefetching backs off under load
instead of adding to it.

\section{Evaluation}
\label{sec:eval}

\paragraph{Setup}
We replay traces against a prototype built on a log-structured store with a
64\,GB DRAM tier and a 2\,TB NVMe tier on a 32-core server. We compare
against LRU-K and ARC without prefetching~\cite{oneil1993lruk,
megiddo2003arc}, ARC with a ghost-list admission filter~\cite{brandao2024ghost},
and ARC with an LSTM prefetcher batched every \SI{200}{\micro\second}, the
best configuration of~\cite{osei2023prefetch}. Traces are a social-graph
store (\emph{Graph}), a session cache (\emph{Session}), an e-commerce
catalogue (\emph{Shop}) and YCSB workload B with Zipfian keys~\cite{cooper2010ycsb}.

\begin{table}[t]
  \centering
  \caption{P99 read latency in microseconds and DRAM hit ratio in percent.
  Lower latency and higher hit ratio are better; the best value in each
  column is set in bold.}
  \label{tab:results}
  \footnotesize
  \setlength{\tabcolsep}{4.2pt}
  \begin{tabular}{@{}l S[table-format=3.0] S[table-format=3.0]
      S[table-format=3.0] S[table-format=3.0] S[table-format=2.1]@{}}
    \toprule
    & \multicolumn{4}{c}{P99 latency (\si{\micro\second})} & {Hit} \\
    \cmidrule(lr){2-5}
    Policy & {Graph} & {Session} & {Shop} & {YCSB-B} & {(\%)} \\
    \midrule
    LRU-K             & 412 & 338 & 501 & 286 & 81.4 \\
    ARC               & 377 & 305 & 466 & 261 & 83.9 \\
    ARC + ghost lists & 351 & 297 & 440 & 255 & 85.2 \\
    ARC + LSTM        & 262 & 248 & 371 & 249 & 88.0 \\
    \midrule
    \sys{}  & \best{201} & \best{209} & \best{318}
                      & \best{246} & \best{91.3} \\
    \bottomrule
  \end{tabular}
\end{table}


\paragraph{Tail latency}
\Cref{tab:results} reports P99 read latency. \sys{} lowers it by
47\% over ARC on \emph{Graph}, where co-access groups are strongest, and by
31\% on \emph{Shop}. On YCSB-B the gain is within noise: keys are drawn
independently, so there are no successors to learn, and the cost gate
correctly declines almost every prefetch. The latency CDF in
\cref{fig:behaviour} shows where the gain comes from: the median is
unchanged, while the fraction of reads served from flash falls from 26\% to
14\%. This is the behaviour we want
from a prefetcher on a workload it cannot help.

\begin{figure}[t]
  \centering
  \begin{tikzpicture}
    \begin{axis}[width=\columnwidth, height=42mm,
        xlabel={DRAM tier as \% of data}, ylabel={hit ratio (\%)},
        xmin=1, xmax=16, ymin=60, ymax=96, xtick={1,2,4,8,16}, xmode=log,
        log basis x=2, xticklabels={1,2,4,8,16},
        legend pos=south east, legend cell align=left]
      \addplot[black!35, dashed] coordinates {(1,62.1) (2,70.4) (4,77.9) (8,83.9) (16,88.6)};
      \addplot[black!60] coordinates {(1,68.3) (2,76.0) (4,82.6) (8,88.0) (16,91.4)};
      \addplot[accent, mark=*, mark size=1.2pt] coordinates {(1,74.8) (2,82.1) (4,87.5) (8,91.3) (16,93.9)};
      \legend{ARC, ARC + LSTM, \sys{}}
    \end{axis}
  \end{tikzpicture}
  \caption{DRAM hit ratio on \emph{Graph} as the fast tier grows. Prefetching
  matters most when DRAM is scarce: at 1\% of the data \sys{}
  matches ARC with four times the memory.}
  \label{fig:capacity}
\end{figure}

\paragraph{Memory pressure}
\Cref{fig:capacity} varies the DRAM tier from 1\% to 16\% of the data. The
benefit of prefetching shrinks as memory grows, because more successors are
already resident, but it never turns negative: the gate keeps wrong guesses
from displacing useful keys. With DRAM at 1\% of the data,
\sys{} reaches the hit ratio that ARC needs 4\% to achieve.

\paragraph{Accuracy and waste}
Of the prefetches \sys{} issues on \emph{Graph}, 73\% are read
within $\Delta$; the LSTM prefetcher, which has no gate, reaches 41\% and
wastes 2.6 times as much flash bandwidth. Raising the threshold $\tau$ in
\eqref{eq:gate} trades coverage for precision smoothly; we found the value
derived from device costs within 5\% of the best setting on every trace
(\cref{fig:behaviour}, centre).

\paragraph{Overhead}
Sketch updates and inference add \SI{2.4}{\micro\second} to the median
request and use 2.7\% of the server's CPU. Retraining takes 11\,s per
shard every five minutes on a background core. Sketches add 410\,B per hot
key, or 3.1\% of the DRAM tier.

\begin{table}[t]
  \centering
  \caption{Ablation on \emph{Graph}: P99 latency, prefetch precision and the
  flash bandwidth spent on prefetches that were never read.}
  \label{tab:ablation}
  \footnotesize
  \begin{tabular}{@{}l S[table-format=3.0] S[table-format=2.0]
      S[table-format=3.0]@{}}
    \toprule
    Variant & {P99 (\si{\micro\second})} & {Prec.\ (\%)} & {Waste (MB/s)} \\
    \midrule
    Full \sys{}             & \best{201} & \best{73} & \best{18} \\
    without cost gate       & 236 & 44 & 71 \\
    without online retraining & 228 & 61 & 29 \\
    sketch size 8 instead of 32 & 219 & 70 & 21 \\
    LSTM instead of trees   & 247 & 69 & 24 \\
    \bottomrule
  \end{tabular}
\end{table}

\paragraph{Ablation}
\Cref{tab:ablation} removes one component at a time. The cost gate matters
most: without it precision drops to 44\% and the wasted bandwidth
quadruples, which in turn slows demand reads that share the device. Online
retraining is the second largest contributor. Replacing the trees by the
LSTM keeps precision almost unchanged but raises latency, because inference
no longer fits on the request path and predictions arrive late.

\paragraph{Drift}
When we switch the \emph{Shop} trace to a new catalogue layout halfway
through, a frozen model loses 60\% of its benefit within an hour, while the
online model recovers its accuracy within two retraining windows
(\cref{fig:behaviour}, right).

\section{Discussion and Limitations}

\paragraph{When prefetching cannot help}
\sys{} learns successors from history, so it helps only where accesses come
in groups that repeat. Uniform or purely Zipfian workloads, like YCSB-B,
offer nothing to learn; the gate makes \sys{} nearly inert on them, but it
still pays the cost of maintaining sketches. A deployment that knows its
workload has no structure should disable it.

\paragraph{Cross-shard successors}
Sketches are local to a shard. When an application reads a key on one shard
and its successor lives on another, \sys{} cannot see the pair. On our
traces 14\% of co-access pairs cross shards; capturing them requires
exchanging sketches between shards, which we leave to future work.

\paragraph{Cost model}
The gate in \eqref{eq:gate} assumes a fixed flash read cost. Under heavy
load the true cost rises with queue depth, and a prefetch that was cheap at
idle becomes expensive. The back-off described in \cref{sec:impl} handles
the extreme case; a gate that reads the current queue depth would handle
the intermediate ones better.

\section{Related Work}

Replacement policies decide what to keep after an
access~\cite{oneil1993lruk,megiddo2003arc}, and learned policies approximate
the optimal offline decision~\cite{song2020lrb}. Admission filters decide
what to let in~\cite{brandao2024ghost}. \sys{} is complementary:
it acts \emph{before} an access, and it reuses the policy's own estimate of
a victim's value in its gate. Learned prefetchers for memory and storage
predict the next address or key with sequence
models~\cite{hashemi2018,osei2023prefetch}; we trade some accuracy for a
model that is cheap enough to run inline and retrain continuously, and we
add the cost model that makes wrong predictions safe. Placement in tiered
stores~\cite{varga2022tiered} decides where a key lives in steady state,
while prefetching responds to the next millisecond.

\section{Conclusion}

Cold reads dominate the tail of tiered key-value stores, and many of them
are predictable from co-access history. \sys{} predicts them with
a small model that runs inline, retrains online, and prefetches only when
the expected saving exceeds the cost of the eviction. It lowers P99 latency
by up to 47\% and never hurts a workload it cannot help. We plan to share
sketches across shards for keys whose successors live elsewhere.

\paragraph{Acknowledgements}
We thank the anonymous reviewers and the operators who shared anonymised
traces. This work was supported by a fictional research grant.

\printbibliography

\end{document}
